\section{Nonlinear Extension and the Path Register}
% -----------------------------------------------------------------------------
\subsection{Exact influence and the Gauss--Newton diagnostic}
Let
\[
L_S(\theta)=L_0(\theta)+\sum_{j\in S}\ell_j(\theta)
\]
and let $\widehat\theta_S$ be a strict local minimum. Write
\[
H_S=\nabla^2L_S(\widehat\theta_S),
\qquad
G_S^{\rm obs}=G_0+\sum_{j\in S}J_j^*W_jJ_j.
\]
The exact infinitesimal response to reweighting point $i$ is
\begin{equation}
v_{i\mid S}^{H}=-H_S^{-1}\alpha_i,
\label{eq:exact-influence}
\end{equation}
whereas the Gauss--Newton observability diagnostic is
\begin{equation}
v_{i\mid S}^{\rm obs}=-(G_S^{\rm obs})^{-1}\alpha_i.
\label{eq:gn-influence}
\end{equation}
They coincide for quadratic models, affine maps under quadratic loss, or when the residual terms in the Hessian are negligible. In general, identifying them would amount to equating Gauss--Newton and the exact Hessian without assumptions \citep{schraudolph2002,martens2020}.

\subsection{Finite insertion through a reweighting path}
Consider
\begin{equation}
L_{S,w}(\theta)=L_S(\theta)+w\ell_i(\theta),
\qquad 0\leq w\leq1.
\end{equation}
Assume there exists a $C^1$ branch of local minima $\gamma_{i\mid S}(w)$ and that the Hessian
\[
H_{S,w}=\nabla^2L_{S,w}(\gamma_{i\mid S}(w))
\]
is invertible along the branch.

\begin{theorem}[Insertion-path equation]
Under the assumptions above,
\begin{equation}
\boxed{
\dot\gamma_{i\mid S}(w)
=-H_{S,w}^{-1}\alpha_i(\gamma_{i\mid S}(w)).
}
\label{eq:insertion-path}
\end{equation}
For any smooth target $\tau$,
\begin{equation}
\boxed{
\tau(\gamma(1))-\tau(\gamma(0))
=-\int_0^1\dd\tau_{\gamma(w)}
\left[H_{S,w}^{-1}\alpha_i(\gamma(w))\right]\dd w.
}
\label{eq:finite-target-effect}
\end{equation}
\end{theorem}

\begin{proof}
The optimality condition is $\dd L_{S,w}(\gamma(w))=0$. Differentiating with respect to $w$ gives
\[
H_{S,w}\dot\gamma(w)+\alpha_i(\gamma(w))=0.
\]
Invertibility of $H_{S,w}$ gives \eqref{eq:insertion-path}, and the integrated chain rule gives \eqref{eq:finite-target-effect}.
\end{proof}

This result is an application of the implicit-function theorem; its role is organizational. It shows that the influence function is the initial velocity of a finite path and that a dynamic value must follow this path. Recent approaches to dynamic valuation and optimal-control sensitivity are therefore direct prior work to discuss, in particular NDDV \citep{liang2024nddv}.

\subsection{Direct contribution and feature-learning feedback}
Fix a connection on parameter space and define along the path
\begin{equation}
G_{S,w}^{\rm obs}(\theta)
=G_0(\theta)+\sum_{j\in S}\Gamma_j(\theta)+w\Gamma_i(\theta).
\end{equation}
Its covariant derivative is
\begin{equation}
\frac{D}{\dd w}G_{S,w}^{\rm obs}(\gamma(w))
=
\underbrace{\Gamma_i(\gamma(w))}_{\text{direct contribution}}
+
\underbrace{\nabla_{\dot\gamma}G_0
+\sum_{j\in S}\nabla_{\dot\gamma}\Gamma_j
+w\nabla_{\dot\gamma}\Gamma_i}_{\text{geometric feedback}}.
\label{eq:geometric-feedback}
\end{equation}
In a linear model with a fixed background metric, the covariant derivatives in the second term vanish. In a network that learns its representations, the displacement induced by the label changes the sensitivities of both present and future points. A frozen regime therefore requires more than a constant $J_i$ for the added point alone: one needs $\nabla G_0=0$, $\nabla J_j=0$, and $\nabla W_j=0$ for every relevant point.

\subsection{Comparing registers over time}
The raw Jacobian $J_t$ is not invariant under reparameterization. To compare functional geometry during learning, fix a declared background metric $g_0$ and a probe set $Z$. We track
\begin{equation}
K_t^{(0)}=J_Z(\theta_t)g_0^\sharp J_Z(\theta_t)^*,
\label{eq:functional-kernel}
\end{equation}
an operator in output space. If $g_0$ transforms tensorially, $K_t^{(0)}$ is invariant under reparameterization. The NTK is the case $g_0=I$ in a fixed parameterization \citep{jacot2018}; the GTK extends this view to other geometries \citep{bai2025gtk}.

The prediction space on the probes is equipped once and for all with its Euclidean inner product. The identification between output vectors and covectors, adjoints, norms, and Hilbert--Schmidt products below are all computed relative to this fixed metric. For two nonzero Hilbert--Schmidt operators $A$ and $B$, we use
\begin{equation}
\operatorname{align}(A,B)
=\frac{\inner{A}{B}_{\rm HS}}{\norm A_{\rm HS}\norm B_{\rm HS}},
\qquad
\operatorname{drift}(A,B)
=\norm{\frac{A}{\norm A_{\rm HS}}-\frac{B}{\norm B_{\rm HS}}}_{\rm HS}
=\sqrt{2\bigl(1-\operatorname{align}(A,B)\bigr)}.
\label{eq:alignment-drift}
\end{equation}
These definitions make the protocol's numerical values reproducible and separate scale variation from change in operator direction. The alignment or normalized drift of $K_t^{(0)}$ is therefore a more defensible diagnostic of feature learning than $\norm{J_t-J_0}$ alone. Dynamic alignment of tangent features has already been studied from other viewpoints \citep{baratin2021}.

The Gram matrix of a hidden representation is also tracked as an auxiliary diagnostic, but it is not a general invariant of the functional model: an invertible transformation compensated in the following layer may change the representation without changing the function. We therefore report it as an architecture- and layer-dependent object, never as an intrinsic invariant.

To compare vectors $v_i(t)$ or $d_i(t)$ belonging to distinct tangent spaces, one must additionally specify a transport: the Levi-Civita connection of $g_0$, transport in a fixed chart, or functional comparison after applying $J_Z$. Without such a convention, a ``vector trajectory of the register'' remains frame-dependent.

\subsection{Dynamic register, spectral measures, and crossings}
At each time $t$, point $i$ has orbit signature
\begin{equation}
\Sigreg_i(t)=\Sigreg\left(G_t^{-1/2}\Gamma_i(t)G_t^{-1/2},
G_t^{-1/2}\alpha_i(t)\right).
\end{equation}
Tracking individually ordered eigenvalues can become unstable at spectral crossings. A continuous representation, symmetric under permutation and adapted to multiplicities, is given by the finite measures
\begin{equation}
\nu_{B_t}=\sum_{\lambda\in\spec(B_t)}m_\lambda(t)\,\delta_\lambda,
\qquad
\mu_{B_t,b_t}=\sum_{\lambda\in\spec(B_t)}\beta_\lambda(t)\,\delta_\lambda.
\label{eq:spectral-measures}
\end{equation}
Their invariant moments are
\begin{equation}
M_k(t)=\int \lambda^k\,\dd\nu_{B_t}(\lambda)=\tr(B_t^k),
\qquad
N_k(t)=\int \lambda^k\,\dd\mu_{B_t,b_t}(\lambda)=b_t^\top B_t^k b_t.
\label{eq:spectral-moments}
\end{equation}
\begin{proposition}[A sufficient number of moments]
Let $B\in\operatorname{Sym}(p)$ and suppose its spectrum contains $r$ distinct eigenvalues. The moments
\[
M_0=p,\;M_1,\ldots,M_p
\]
determine the characteristic polynomial of $B$ through Newton's identities, and hence the spectral measure $\nu_B$ with multiplicities. Once the distinct eigenvalues $\lambda_1,\ldots,\lambda_r$ are known, the moments
\[
N_0,\ldots,N_{r-1}
\]
determine the energies $\beta_{\lambda_j}$ through the Vandermonde system
\[
N_k=\sum_{j=1}^r\beta_{\lambda_j}\lambda_j^k,
\qquad 0\le k\le r-1.
\]
\end{proposition}
\begin{proof}
The power sums $M_k=\tr(B^k)$ recursively determine the elementary coefficients of the characteristic polynomial. The second system is linear, with matrix $(\lambda_j^k)_{0\le k<r,1\le j\le r}$, which is invertible when the $\lambda_j$ are distinct.
\end{proof}
These formulas are exact but may be numerically ill-conditioned when eigenvalues are close, because the Vandermonde matrix becomes unstable. For dynamic tracking, however, the moments $M_k(t)$ and $N_k(t)$ remain continuous in the coefficients of $(B_t,b_t)$ even when eigenvalue labels swap. The discrete signature is useful for instantaneous interpretation; measures, better-conditioned polynomial bases, or resolvents are preferable for trajectories.

The trajectory may reveal transitions absent from the final state: an informative point becomes redundant, a previously invisible direction becomes observable, or a coherent point becomes conflicting after representation reorganization. This dynamic accounting should not be confused with a claim of priority over all temporal attribution; it provides a particular tensorial decomposition of the mechanisms being tracked.

% -----------------------------------------------------------------------------
